% Part: lambda-calculus
% Chapter: lambda-definability
% Section: lists

\documentclass[../../../include/open-logic-section]{subfiles}

\begin{document}

% SEGMENT OLP-0650-S01
\olfileid{lam}{ldf}{lst}
\olsection{பட்டியல்கள்}

$[M_1, M_2, \ldots, M_n]$ என்ற பட்டியலைப் பின்வருமாறு வரையறுக்கலாம்:
\[
  [M_1, M_2, \ldots, M_n] \ident \lambd[sz][s M_1 (s M_2 (\ldots (s M_n z)))]
\]
முறைசாரா வகையில் சொன்னால், பட்டியலைக் குறிக்கும் சொல் அதன்
‘திரட்டுச் சார்பு’ ஆகும்: அது இரண்டு சொற்களை ஏற்கும் சார்பு.
முதலாவது, புதிய உறுப்பையும் இதுவரை திரட்டிய மதிப்பையும் பெற்று
புதிய திரண்ட மதிப்பைத் தரும் சார்பு; இரண்டாவது, தொடக்கத் திரண்ட
மதிப்பு. இவ்வாறு உறுப்புகளை ஒவ்வொன்றாகத் திரட்டி இறுதியில்
விளைவைத் தருகிறது.

% SEGMENT OLP-0650-S02
\begin{digress}
  சார்புசார் நிரலாக்கத்தை அறிந்தவர்களுக்கு, இந்த வரையறை
  \emph{வலப்புற மடிப்பு} என்பதற்கு ஒத்ததாகத் தோன்றும்:
  ஒரு பட்டியல், அதிலுள்ள உறுப்புகள்மீது வலப்புற மடிப்பைச் செய்யும்
  சார்பாக வரையறுக்கப்படுகிறது.
\end{digress}

% SEGMENT OLP-0650-S03
இந்தக் குறியாக்கத்தின் அடிப்படையில் பயனுள்ள சில சார்புகளை எளிதில்
வரையறுக்கலாம்:
\begin{align*}
  \fn{Sum} & \ident
  \lambd[l][l \, (\lambd[xa][\fn{Add} \, x \, a])\,  \num{0}]\\
  \fn{Len} & \ident
  \lambd[l][l \, (\lambd[xa][\fn{Add} \, a \, \num{1}]) \num{0}]
\end{align*}

$\fn{Sum}$ ஆனது சர்ச் எண்ணுருக்களின் பட்டியலில் உள்ளவற்றின்
கூட்டுத்தொகையைக் கணக்கிடுகிறது. தொடக்க மதிப்பு $\num{0}$
ஆக இருக்கும் திரட்டலைப் பட்டியல்மீது செய்கிறது; ஒவ்வோர் உறுப்புக்கும்,
புதிய மதிப்பாக $\fn{Add} \, x\, a$ ஐத் தருகிறது
(இங்கு $x$ தற்போதைய உறுப்பு, $a$ இதுவரை திரண்ட மதிப்பு).
விளைவு எல்லா உறுப்புகளின் கூட்டுத்தொகை.

% SEGMENT OLP-0650-S04
\begin{prob}
  $\fn{Len}$ என்ன செய்கிறது? விளக்குக.
\end{prob}

% SEGMENT OLP-0650-S05
\begin{prob}
  ஒரு பட்டியலைத் தலைகீழாக்கும் சார்பை வரையறுக்க.
\end{prob}
\end{document}
